\Needspace{12\baselineskip}
\originalchapter{补充公开实验：从粒子到量子操控}
以下主题均有已取得的Junior Lab官方手册，而非只见目录。它们扩展此前九个测量模块的覆盖，沿装置、物理模型、分析和误差的相同论证链展开。实体装置并未在本项目运行，涉及低温、放射源或激光的功能说明不能代替实验室操作规程；数值例子按模拟/模型明确标记。

\originalsection{宇宙线muon的速度与寿命}
速度测量用上下两块闪烁体、光电探测器、判别/符合电路和时间测量单元。相同粒子跨已知间距 $L$，到达时间差 $\Delta t$ 给 $v=L/\Delta t$。两个通道的电缆延迟、判别阈值和脉冲上升时间会造成固定偏移及time walk；先以同一脉冲或改变几何的相反顺序标定电子学延迟，再解释纳秒量级飞行时间。

有限面积接受倾斜轨迹，实际路长不总等于垂直间距，若不限定角度会把速度低估。测量 $v$ 接近 $c$ 时，误差不能支持“超过光速”便丢点；应以完整几何和时间分辨模型判断。不同 $L$ 对平均时间作直线拟合，截距可暴露共同电子学延迟，斜率更接近 $1/v$。

寿命测量让muon停在厚探测介质，一次脉冲为停入，延迟脉冲为衰变产物。若独立衰变率为常 $1/\tau$，生存概率 $P(t)=\e^{-t/\tau}$，等待密度 $p=\tau^{-1}\e^{-t/\tau}$。偶然背景近常数，仪器盲区和触发窗改变可见分布；拟合 $A\e^{-t/\tau}+B$ 或有限窗条件似然，不能任意从零时刻取对数拟合。

负muon可被原子核俘获，使某些材料中的有效消失率 $1/\tau_{
m eff}=1/\tau_{\rm decay}+\lambda_{
m capture}$，混合正负muon可能不是单一指数。讨论宇宙线从大气到地面的存活还需飞行高度和Lorentz因子 $\gamma$，实验室寿命与运动寿命 $\gamma\tau_0$ 区分。本卷没有凭停止衰变曲线直接推算未给的初始天空通量。

\originalsection{Rutherford散射的角分布和薄靶}
非相对论alpha粒子受中心库仑势，能量和角动量守恒给冲量参数 $b=[Z_1Z_2e^2/(8\pi\epsilon_0E)]\cot(\theta/2)$。由轴对称面积 $2\pi b\dd b$ 对立体角 $2\pi\sin\theta\dd\theta$ 的比，
\[
\frac{\dd\sigma}{\dd\Omega}=\left(\frac{Z_1Z_2e^2}{16\pi\epsilon_0E}\right)^2\frac1{\sin^4(\theta/2)}.
\]
该推导假定单散射、未接近核力尺度、电子屏蔽对所考察角度不主导。$E^{-2}$、$Z_2^2$ 和强角依赖是可独立比较的预测。

装置用准直源、薄金属靶、可转角半导体探测器和能谱计数。期望计数 $n=\Phi t\,n_t\Delta\Omega\,\epsilon_{
m det}\,\dd\sigma/\dd\Omega$，$n_t$ 是单位面积靶核数，$\Delta\Omega$ 需按真实孔径/距离积分，$\Phi$ 是入射率。源率与靶厚不明时可比较相对角分布，但不能报告绝对截面。

薄靶仍有能量损失及多重散射，角接受宽需卷积Rutherford陡曲线。向前小角时有限角宽尤其影响平均。背景和未散射束要独立控制；改变靶厚若计数不线性或峰明显展宽，应检查单散射假设。源手册图和第三方实验曲线没有复制。

\originalsection{氢样光谱：波长校准与精细结构}
扫描单色仪用狭缝、准直、色散元件和探测器将波长变为位置/读数。汞谱线先标定波长关系，狭缝宽改变分辨和强度。氢Balmer线满足 $1/\lambda=R_H(1/2^2-1/n^2)$，所以绘制倒波长对 $1/4-1/n^2$，斜率给Rydberg常数。波长误差与其倒数误差不同，校准参数在各谱线中共享。

从库仑束缚能 $E_n=-\mu_r e^4/[2(4\pi\epsilon_0)^2\hbar^2n^2]$ 可见约化质量进入；核无限重只是近似。氘与氢的同位素位移因此不靠改变电子电荷解释。钠一个价电子在闭壳层外，但芯屏蔽和量子缺陷使其不完全是纯氢公式；精细结构与自旋轨道耦合又产生双线。

拟合谱线应包含仪器线形与背景，分辨不了的双线可能只给加权峰心。用尖峰最大读数作“真实谱线”低估不确定度。根据校准残差、不同狭缝和重复扫描比较，才能区分物理分裂与仪器漂移。

\originalsection{X射线：连续谱、特征线和吸收}
电子经电压 $V$ 加速后撞靶，最大可给光子能量 $eV$，连续轫致谱截止 $\lambda_{min}=hc/(eV)$。靶内壳空位由外层电子填入产生特征X射线，与电压连续谱机制不同；电压超过激发阈值才有相应线。简单Moseley标度 $\nu\propto(Z-\sigma_{
m screen})^2$ 的有效屏蔽参数不能对所有壳层通用。

固态能量探测器读脉冲高度，经已知谱线标定后得到谱。吸收近似 $I(E,x)=I_0(E)\e^{-\mu(E)x}$，只有窄单能线才可用一个 $\mu$；多能束穿过材料变硬，平均有效吸收随厚度变。吸收边对应内壳阈值，不能拿连续光滑幂律覆盖全部能区。

放射荧光分析需入射谱、几何、自吸收和探测效率。Compton散射、逃逸峰、电子学pileup可能形成次峰：双脉冲重叠的总能峰随计数率变化，不是新元素线。死时间与能谱标定互相连接，补充手册因此与第3章的测量链共用误差分析。

\originalsection{铷光学泵浦和磁共振}
装置用铷蒸气池、D1光、圆偏振光学、三轴磁场线圈与RF场。圆偏振选择跃迁 $\Delta m=+1$ 或 $-1$，多次吸收/发射把群体积累到不再吸收该偏振的端态，透射增加。这是人口重新分布而非简单加热增强磁性。

弱场超精细子空间Zeeman能 $E_m\simeq g_F\mu_BmB$，相邻共振频率 $\nu=|g_F|\mu_BB/h$。两个天然同位素具有不同核自旋与 $g_F$，因而共振斜率不同。实际地磁与线圈场向量相加，扫描一线圈的共振频率不是自动对电流直线，除非其他分量已补偿。

改变偏振、断光/恢复光和RF驱动，可测人口弛豫与泵浦时间。光强、碰撞、温度和磁场不均匀改变线宽；功率展宽与真实弛豫不易只靠一条曲线分开。共振截距可用于地磁估计，仍需线圈常数、方向和频率校准。泵浦装置工作解释与第4章热平衡NMR初态不同，二者不能混用同一人口模型。

\originalsection{M\"ossbauer：无反冲谱和超精细分裂}
自由核发/吸光子有反冲能 $E_R=E_\gamma^2/(2Mc^2)$，窄自然线时发射与吸收能不重合。晶格可整体承受动量，部分事件无声子激发、近无反冲，才产生极窄共振。Debye--Waller型无反冲比例取决于原子位移涨落与光子波数，不是每次发射都无反冲。

手册以机械驱动的源相对吸收体速度 $v$ 扫描，Doppler能移 $\Delta E\simeq E_\gamma v/c$ 将毫米/秒转换为极小能量差。铁57约14 keV线的超精细磁分裂、四极分裂及化学位移分别反映磁场、电场梯度及核电子密度环境。峰数/强度受选择定则与几何极化影响，不能单从“六峰”唯一推原子所有参数。

驱动器速度需标定与线性检查，源/吸收体厚度可引入饱和展宽，常数背景和机械零点进拟合。Lorentz线型宽度与峰心统计误差不同。化学位移应给相对参考材料和温度，第二阶Doppler温度项也可能影响中心；不把任意位移都归于化学价态变化。

\originalsection{超导、Meissner与Josephson测量}
手册包含液氦低温装置、磁通排斥、持久流和Josephson现象。电阻为零与Meissner排斥是不同观测：理想零电阻导体能冻结已有磁通，超导平衡态则在适用临界条件下排斥。London模型给 $\nabla^2\mathbf B=\mathbf B/\lambda_L^2$，半空间场随深度 $\e^{-z/\lambda_L}$ 衰减；实际样品形状的退磁场会改变外测线圈信号。

Type-I与Type-II的临界场结构不同，后者在 $H_{c1}<H<H_{c2}$ 有涡旋混合态。零场冷却与有场冷却、磁通钉扎和滞后需分别记录；若只看一次线圈电压跳变，尚不能辨认全部平衡磁性。

Josephson结两端相位差 $\phi$ 满足 $I_s=I_c\sin\phi$、$\dot\phi=2eV/\hbar$；恒电压给频率 $\nu=2eV/h$。在外微波频率 $\nu_d$ 下可有Shapiro台阶 $V_n=nh\nu_d/(2e)$。SQUID两结的相位受磁通约束，临界流随磁通周期为 $\Phi_0=h/(2e)$。引线电阻、寄生磁场、热循环和电磁屏蔽都影响可观测曲线；本卷未伪造低温实测台阶。

\originalsection{无Doppler谱与交叉共振}
热原子沿束方向速度 $v$，一束的共振条件带 $kv$，温度速度分布给DopplerGaussian展宽。强泵浦和弱探测反向传播，二者同时与同一跃迁共振只在近零轴向速度，强束饱和人口使探测吸收出现窄凹陷。该饱和吸收消去一阶速度展宽，但自然线宽、碰撞和功率展宽仍存在。

两个相近跃迁也可产生交叉共振：某速度组对一束共振于第一线、另一束于第二线，激光频率位于两线中点。交叉峰不是新的原子能级。扫描激光电流/压电非线性，需参考腔或已知线间距标频轴，重复扫向检查滞后。

光强过高、泵探几何失配及剩余磁场会改变窄线，频率调制/锁相又输出线形导数而非原吸收。拟合时必须使用与检测方式一致的模型，不能把导数零点两侧峰间距误当两个物理跃迁。

\originalsection{NMR量子信息的态、脉冲和读出}
手册用氯仿的质子与碳13两自旋，旋转系近似Hamilton量 $H/\hbar=\Delta\omega_HI_z^H+\Delta\omega_CI_z^C+2\pi JI_z^HI_z^C$。$I_z=\sigma_z/2$，演化 $U=\e^{-iHt/\hbar}$；RF脉冲给单比特旋转 $R_x(\theta)=\e^{-i\theta\sigma_x/2}$。耦合演化给条件相位，经单自旋相位和基变换可构成CNOT。

热平衡NMR样品是混态，近高温 $\rho\simeq I/4+\epsilon\Delta\rho$，信号来自微小偏离部分。伪纯态的计算信号可模拟纯态算法，但不等于已准备大比例真正纯态。读取的是整体横磁化，接收复信号 $s(t)=\operatorname{tr}[\rho(t)(I_x+iI_y)]$，频谱峰幅与特定密度矩阵非对角元有关；不能把单个峰直接当全部计算基概率。

作为实际运行的教学验证，附带NumPy脚本构造两比特Hadamard、四种函数oracle、CNOT与Deutsch--Jozsa电路。先置 $|0,1\rangle$，两比特Hadamard后经 $U_f|x,y\rangle=|x,y\oplus f(x)\rangle$ 再对第一比特Hadamard：常数函数测第一比特0的概率为1，平衡函数为0。矩阵幺正与四种输出实际核验，但没有假装在NMR实体硬件执行算法。

\originalsection{Raman项目资料的边界与转动谱}
公开Raman文件是Hengyun Zhou、Francisco Machado的2015学生项目提案，性质与标准完成实验手册不同。本卷保留这一署名与角色，采用独立的物理解释，不把提案的目标写成已经取得的实验结果。

入射光经虚中间态散射，弹性Rayleigh线频率不变，Stokes/anti-Stokes线差给分子能差；热平衡强度比受Boltzmann人口及频率因子控制。二原子刚转子 $E_J=BJ(J+1)$，$B=\hbar^2/(2I)$，转动Raman常有 $\Delta J=\pm2$，Stokes能移 $\Delta E=B(4J+6)$，相邻线间距 $4B$。由 $I=\mu_r r_e^2$ 可反推键长，但需波数标定与精细模型。

同核分子的核自旋统计限制偶/奇 $J$ 人口，可能造成交替线强。实际提案的O2/N2谱还受转动、电子自旋和仪器线宽影响。高强Rayleigh杂散光会掩盖弱Raman线，光路滤除、波数校准与相机背景构成主要测量链。这里未生成假称原项目已测得的分子谱。
